Táto práca overila, že vybrané kombinatorické úlohy možno spracovať cez spoločný redukčný postup. Rozhodovacie premenné, ich konečné domény a obmedzenia sa najprv zapíšu ako problém splniteľnosti obmedzení. Z neho sa potom vytvorí Boolova formula a jej \acrshort{knf} kódovanie vhodné pre \acrshort{sat} solver. Pri tomto postupe nie je potrebné zachovať všetky modely pôvodného zápisu. Podstatné je zachovať riešiteľnosť, preto má v práci dôležité miesto ekviosplniteľnosť a Tseitinova transformácia.

Teoretická časť zaviedla pojmy splniteľnosti, modelu, logickej ekvivalencie, problému splniteľnosti obmedzení a redukcie na \acrshort{sat}. Na tomto základe boli formulované grafové úlohy farbenia vrcholov, farbenia hrán, nezávislej množiny a snarkov. Implementácia zahŕňa tie typy vstupov, ktoré aplikácia priamo podporuje. Čiže Boolove formuly, priame \acrshort{sat} inštancie, všeobecné \acrshort{csp} inštancie, regulárne farbenie vrcholov a hrán.

Výsledkom praktickej časti je webová aplikácia, v ktorej používateľ zadáva problém v doménovom tvare a systém sa stará o jeho prevod na \acrshort{sat}, spustenie externého solvera a interpretáciu získaného modelu. Návrh oddeľuje doménové pravidlá od webového rozhrania, databázy a technického spúšťania solverov. Riešenie úloh prebieha asynchrónne cez samostatného pracovníka, takže používateľská požiadavka iba zaradí úlohu na spracovanie a výsledok sa uloží po skončení výpočtu.

Testovacia časť zaznamenáva, že pri menších vstupoch je možné prejsť celým tokom od zadania problému až po zobrazenie výsledku. Zároveň však ukázala hranicu aktuálneho návrhu. Plánované výkonnostné merania nebolo možné uzavrieť ako spoľahlivé porovnanie solverov CaDiCaL a CryptoMiniSat, pretože pri väčších grafových vstupoch zlyhávalo už spracovanie dát v aplikácii. Namerané časy by preto nevyjadrovali iba správanie solverov.

Hranicou aktuálnej verzie je najmä práca s veľkými textovými artefaktmi, ktoré sa ukladajú priamo v databázovom riadku. Ide najmä o serializovanú inštanciu a vygenerované \acrshort{dimacs} kódovanie. Ďalší vývoj by mal tieto dáta oddeliť od metadát problému, napríklad uložením do samostatných súborov alebo objektového úložiska. Po tejto úprave bude možné zopakovať výkonnostné testy na väčších grafoch, porovnať použité solvery a rozšíriť implementáciu o ďalšie úlohy opísané v teoretickej časti.
